\documentclass[11pt]{article}

%%%%%%%%%%%%% Package imports %%%%%%%%%%%%%
% Uses Latin Modern + T1 when available; falls back to default fonts otherwise.
\IfFileExists{lmodern.sty}%
  {\usepackage[T1]{fontenc}\usepackage{lmodern}\usepackage{microtype}}%
  {\usepackage[T1]{fontenc}\usepackage[expansion=false]{microtype}}
\usepackage[backend=biber, style=numeric]{biblatex}
\usepackage{enumitem}
\usepackage[margin=1in]{geometry}
\usepackage{amsmath}
\usepackage{amssymb}
\usepackage{bm}
\usepackage{booktabs}
\usepackage{algorithm}
\usepackage{algpseudocode}
\usepackage[hidelinks]{hyperref}

\addbibresource{iekf_notes.bib}

%%%%%%%%%%%%% Macro definitions %%%%%%%%%%%%%
\setlist{itemsep=2pt,topsep=4pt}

\newcommand{\x}{\bm{x}}
\newcommand{\y}{\bm{y}}
\newcommand{\uvec}{\bm{u}}
\newcommand{\w}{\bm{w}}
\newcommand{\vvec}{\bm{v}}
\newcommand{\dyn}{\bm{f}}
\newcommand{\meas}{\bm{h}}
\newcommand{\F}{\bm{F}}
\newcommand{\Hm}{\bm{H}}
\newcommand{\K}{\bm{K}}
\newcommand{\Pm}{\bm{P}}
\newcommand{\Q}{\bm{Q}}
\newcommand{\R}{\bm{R}}
\newcommand{\Sm}{\bm{S}}
\newcommand{\I}{\bm{I}}
\newcommand{\xbar}{\bar{\x}}
\newcommand{\xhat}{\hat{\x}}
\newcommand{\Pbar}{\bar{\Pm}}
\newcommand{\Phat}{\hat{\Pm}}
\newcommand{\inn}{\bm{\nu}}
\newcommand{\Normal}{\mathcal{N}}
\newcommand{\iter}[1]{^{(#1)}}

%%%%%%%%%%%%% Title card content %%%%%%%%%%%%%
\title{The Iterated Extended Kalman Filter:\\
Equations, Trade-offs, Diagnostics, and Use on Real Data}
\author{Andrew\\ \small Director of Research \& Development, Terminus Stratospherics}
\date{\today}

%%%%%%%%%%%%% Begin document %%%%%%%%%%%%%
\begin{document}
\maketitle
\tableofcontents

%======================================================================
\section{Model Setup}

Consider the discrete-time nonlinear system with additive Gaussian noise
\begin{align}
  \x_k &= \dyn(\x_{k-1}, \uvec_{k-1}) + \w_{k-1}, & \w_{k-1} &\sim \Normal(\bm{0}, \Q_{k-1}), \label{eq:dyn}\\
  \y_k &= \meas(\x_k) + \vvec_k, & \vvec_k &\sim \Normal(\bm{0}, \R_k), \label{eq:meas}
\end{align}
with $\x \in \mathbb{R}^n$ and $\y \in \mathbb{R}^m$. If the dynamics are continuous,
$\dyn$ denotes the integrated flow over $[t_{k-1}, t_k]$ and the Jacobian $\F$ below is the
state transition matrix $\boldsymbol{\Phi}(t_k, t_{k-1})$, obtained by integrating
$\dot{\boldsymbol{\Phi}} = \mathbf{A}(t)\boldsymbol{\Phi}$, $\boldsymbol{\Phi}(t_{k-1},t_{k-1}) = \I$,
alongside the state.

%======================================================================
\section{Prediction (Time Update)}

The IEKF prediction step is identical to that of the EKF.

\paragraph{State propagation.}
\begin{equation}
  \xbar_k = \dyn(\xhat_{k-1}, \uvec_{k-1}).
\end{equation}

\paragraph{Dynamics Jacobian.}
\begin{equation}
  \F_{k-1} = \left.\frac{\partial \dyn}{\partial \x}\right|_{\xhat_{k-1}}.
\end{equation}

\paragraph{Covariance propagation.}
\begin{equation}
  \Pbar_k = \F_{k-1}\Phat_{k-1}\F_{k-1}^\top + \Q_{k-1}.
\end{equation}

The standard IEKF does \emph{not} iterate the prediction step; see
Section~\ref{sec:variants} for variants that do.

%======================================================================
\section{Iterated Measurement Update}

\subsection{Origin: MAP Estimation via Gauss--Newton}

The EKF update linearizes $\meas$ once, about the prior $\xbar_k$. The IEKF instead seeks the
maximum a posteriori (MAP) estimate under the Gaussian prior and Gaussian likelihood, i.e.\ the
minimizer of
\begin{equation}
  J(\x) = \tfrac{1}{2}\bigl(\y_k - \meas(\x)\bigr)^\top \R_k^{-1}\bigl(\y_k - \meas(\x)\bigr)
        + \tfrac{1}{2}(\x - \xbar_k)^\top \Pbar_k^{-1}(\x - \xbar_k).
  \label{eq:cost}
\end{equation}
Bell and Cathey~\cite{bell1993} showed that the IEKF iteration is exactly the
Gauss--Newton method applied to~\eqref{eq:cost}: the filter relinearizes $\meas$ about
successively better estimates of the posterior mode.

\subsection{The Iteration}

Initialize at the prior,
\begin{equation}
  \x\iter{0} = \xbar_k ,
\end{equation}
and for $i = 0, 1, 2, \dots$ perform the following.

\begin{enumerate}[label=\textbf{Step \arabic*.}, leftmargin=*, itemsep=1ex]
  \item \textbf{Relinearize} the measurement model at the current iterate:
  \begin{equation}
    \Hm\iter{i} = \left.\frac{\partial \meas}{\partial \x}\right|_{\x\iter{i}}.
  \end{equation}

  \item \textbf{Innovation covariance} (the prior covariance is fixed; only $\Hm$ changes):
  \begin{equation}
    \Sm\iter{i} = \Hm\iter{i}\Pbar_k\Hm^{(i)\top} + \R_k .
  \end{equation}

  \item \textbf{Gain:}
  \begin{equation}
    \K\iter{i} = \Pbar_k\Hm^{(i)\top}\bigl(\Sm\iter{i}\bigr)^{-1}.
  \end{equation}

  \item \textbf{State iterate} (the key equation):
  \begin{equation}
    \boxed{\;
      \x\iter{i+1} = \xbar_k + \K\iter{i}\Bigl[\y_k - \meas\bigl(\x\iter{i}\bigr)
                     - \Hm\iter{i}\bigl(\xbar_k - \x\iter{i}\bigr)\Bigr]
    \;}
    \label{eq:iekf-update}
  \end{equation}
  The correction term $-\Hm\iter{i}(\xbar_k - \x\iter{i})$ is what makes this Gauss--Newton
  rather than a naive re-application of the EKF update: every iterate is anchored to the prior
  $\xbar_k$, not to the previous iterate. For $i = 0$ the term vanishes and the EKF update is
  recovered exactly.

  \item \textbf{Convergence check:} stop when
  \begin{equation}
    \bigl\|\x\iter{i+1} - \x\iter{i}\bigr\| < \epsilon
    \quad\text{or}\quad i + 1 = i_{\max}.
  \end{equation}
  A scale-aware alternative is the Mahalanobis step size
  $\Delta^\top \Pbar_k^{-1} \Delta < \epsilon$ with $\Delta = \x\iter{i+1} - \x\iter{i}$,
  which avoids mixing units (e.g.\ km and rad) in a single norm.
\end{enumerate}

\subsection{Final Posterior}

Let $N$ denote the converged iterate. Then
\begin{equation}
  \xhat_k = \x\iter{N}.
\end{equation}
Recompute $\Hm$ and $\K$ at $\x\iter{N}$ and update the covariance in Joseph form for
numerical robustness:
\begin{equation}
  \Phat_k = (\I - \K\Hm)\,\Pbar_k\,(\I - \K\Hm)^\top + \K\R_k\K^\top .
  \label{eq:joseph}
\end{equation}
This is the inverse of the Gauss--Newton Hessian approximation of $J$ at the mode, i.e.\ a
Laplace-type approximation of the posterior covariance.

\begin{algorithm}[htbp]
\caption{IEKF measurement update (with optional damping)}
\label{alg:iekf}
\begin{algorithmic}[1]
\Require Prior $\xbar_k$, $\Pbar_k$; measurement $\y_k$; $\R_k$; tolerance $\epsilon$; $i_{\max}$
\State $\x\iter{0} \gets \xbar_k$
\For{$i = 0, 1, \dots, i_{\max}-1$}
  \State $\Hm\iter{i} \gets \partial\meas/\partial\x\,|_{\x\iter{i}}$
  \State $\Sm\iter{i} \gets \Hm\iter{i}\Pbar_k\Hm^{(i)\top} + \R_k$
  \State $\K\iter{i} \gets \Pbar_k\Hm^{(i)\top}(\Sm\iter{i})^{-1}$
  \State $\x_{\mathrm{GN}} \gets \xbar_k + \K\iter{i}\bigl[\y_k - \meas(\x\iter{i}) - \Hm\iter{i}(\xbar_k - \x\iter{i})\bigr]$
  \State $\Delta \gets \x_{\mathrm{GN}} - \x\iter{i}$
  \State $\alpha \gets 1$
  \While{$J(\x\iter{i} + \alpha\Delta) > J(\x\iter{i})$ \textbf{and} $\alpha > \alpha_{\min}$} \Comment{optional backtracking}
    \State $\alpha \gets \alpha/2$
  \EndWhile
  \State $\x\iter{i+1} \gets \x\iter{i} + \alpha\Delta$
  \If{$\Delta^\top\Pbar_k^{-1}\Delta < \epsilon$} \textbf{break} \EndIf
\EndFor
\State $\xhat_k \gets \x\iter{i+1}$; recompute $\Hm$, $\K$ at $\xhat_k$
\State $\Phat_k \gets (\I - \K\Hm)\Pbar_k(\I - \K\Hm)^\top + \K\R_k\K^\top$
\end{algorithmic}
\end{algorithm}

\subsection{Damped / Line-Search IEKF}

Pure Gauss--Newton can overshoot or diverge when the prior is poor or $\meas$ is strongly
curved. The damped IEKF~\cite{skoglund2015} computes the full Gauss--Newton step
\begin{equation}
  \Delta\iter{i} = \x_{\mathrm{GN}}\iter{i+1} - \x\iter{i}
\end{equation}
and takes a partial step
\begin{equation}
  \x\iter{i+1} = \x\iter{i} + \alpha_i\,\Delta\iter{i}, \qquad \alpha_i \in (0, 1],
\end{equation}
with $\alpha_i$ chosen by backtracking until $J$ decreases. A Levenberg--Marquardt variant,
which inflates the prior-information term, achieves the same effect. The modification is cheap
and turns the IEKF from ``usually better'' into ``reliably no worse'' than the EKF.

\subsection{Attitude and Other Manifold States}

For states containing quaternions or rotation matrices, iterate in the \emph{error state}.
The prior term becomes the manifold difference $\x\iter{i} \boxminus \xbar_k$, so in
\eqref{eq:iekf-update} the quantity $(\xbar_k - \x\iter{i})$ is replaced by
$-(\x\iter{i} \boxminus \xbar_k)$, and each update is retracted with $\boxplus$.
Strictly, $\Pbar_k$ should also be transported to the tangent space at the new iterate.
This correction is often small but matters for large attitude corrections; in coupled
orbit--attitude problems, handling it properly is typically more important than the
iteration count.

%======================================================================
\section{Related Variants}
\label{sec:variants}

\begin{table}[htbp]
\centering
\begin{tabular}{@{}lp{0.5\textwidth}@{}}
\toprule
\textbf{Variant} & \textbf{What is iterated} \\
\midrule
IEKF (this document) & Measurement linearization only \\
Iterated EKS / Gauss--Newton smoother & Whole trajectory, both $\dyn$ and $\meas$ (batch MAP) \\
Iterated UKF / IPLF & Statistical (sigma-point) linearization of $\meas$ about the current posterior \\
Damped IEKF & Measurement linearization with line search or LM damping \\
\bottomrule
\end{tabular}
\caption{IEKF and related iterated estimators.}
\end{table}

%======================================================================
\section{Advantages and Disadvantages}

\subsection{Advantages}
\begin{itemize}
  \item \textbf{Better handling of measurement nonlinearity} than the EKF (range/bearing,
        angles-only, star-tracker/line-of-sight, Doppler). Reduces linearization bias in the
        update.
  \item \textbf{Clear optimization interpretation.} Each update is the MAP estimate given the
        Gaussian prior, which provides principled tools: cost monotonicity checks, damping, and
        convergence tests.
  \item \textbf{Same ingredients as the EKF.} No new Jacobians, no sigma points; usually a
        small code change.
  \item \textbf{Adaptive cost.} Mildly nonlinear updates converge in one or two iterations, so
        extra computation is spent only where nonlinearity matters.
  \item \textbf{Robust to low-to-high information transitions}, such as the first good
        measurement after a long coast, when the prior is large relative to $\R$---precisely
        where the EKF tends to diverge.
\end{itemize}

\subsection{Disadvantages}
\begin{itemize}
  \item \textbf{No help with dynamics nonlinearity.} Prediction remains a first-order
        covariance propagation; long propagation arcs or strongly nonlinear dynamics can still
        corrupt $\Pbar_k$.
  \item \textbf{Higher, variable cost} per update (multiple evaluations of $\meas$, $\Hm$, and
        $\Sm^{-1}$). Awkward under hard real-time budgets; cap $i_{\max}$.
  \item \textbf{Undamped Gauss--Newton can diverge or oscillate} when the prior is badly
        wrong, potentially performing worse than the EKF.
  \item \textbf{Covariance is still a local linearization}, and can be overconfident when the
        posterior is non-Gaussian, skewed, or banana-shaped (e.g.\ angles-only orbit
        determination).
  \item \textbf{Unimodal.} The filter converges to \emph{a} mode and cannot represent
        ambiguity (multiple solutions, data-association uncertainty).
  \item \textbf{Little benefit when $\R$ dominates} $\Hm\Pbar\Hm^\top$: the update is small,
        linearization error negligible, and the extra computation is wasted.
  \item \textbf{Sensitive to errors in $\meas$.} Iterating toward a biased measurement model fits
        the bias more confidently.
\end{itemize}

%======================================================================
\section{Diagnostics for Filter Health}

\subsection{Innovation-Based Tests (No Truth Required)}

The innovation (pre-fit residual) and its predicted covariance, evaluated at the prior, are
\begin{equation}
  \inn_k = \y_k - \meas(\xbar_k), \qquad
  \Sm_k = \Hm_k\Pbar_k\Hm_k^\top + \R_k .
\end{equation}
For a healthy filter, $\inn_k$ is zero-mean, white, and has covariance $\Sm_k$. Three tests
check these properties.

\paragraph{Normalized Innovation Squared (NIS).}
\begin{equation}
  \epsilon_{\nu,k} = \inn_k^\top \Sm_k^{-1}\inn_k \sim \chi^2_m .
\end{equation}
Over a window of $N$ steps, $\sum_{k} \epsilon_{\nu,k} \sim \chi^2_{Nm}$, so the time-averaged
NIS should lie within the two-sided 95\% bounds
\begin{equation}
  \frac{1}{N}\Bigl[\chi^2_{Nm}(0.025),\; \chi^2_{Nm}(0.975)\Bigr].
\end{equation}
\begin{itemize}
  \item NIS persistently too high: the filter is overconfident ($\Q$ or $\R$ too small, or the
        model is wrong).
  \item NIS persistently too low: the filter is too conservative.
\end{itemize}

\paragraph{Whiteness (autocorrelation) test.}
For the normalized innovation sequence, compute the sample autocorrelation $\hat\rho(j)$ for
lags $j \ge 1$. Approximately 95\% of values should satisfy
\begin{equation}
  |\hat\rho(j)| \le \frac{1.96}{\sqrt{N}} .
\end{equation}
Correlated innovations indicate unmodeled dynamics, colored noise, or a mistuned $\Q$.

\paragraph{Bias test.}
The sample mean of each normalized innovation component should be approximately zero. A
persistent offset suggests a sensor bias that should be estimated as an additional state.

\subsection{Truth-Based Tests (Simulation and Monte Carlo)}

\paragraph{Normalized Estimation Error Squared (NEES).}
\begin{equation}
  \epsilon_{x,k} = (\x_k - \xhat_k)^\top \Phat_k^{-1} (\x_k - \xhat_k) \sim \chi^2_n .
\end{equation}
Averaging over $M$ Monte Carlo runs gives the ANEES, which is checked against
$\chi^2_{Mn}/M$ bounds. This is the gold-standard consistency test~\cite{barshalom2001}.

\paragraph{Error versus $\pm 3\sigma$ envelopes.}
Plot each state's estimation error against $\pm 3\sigma_j$, $\sigma_j = \sqrt{\hat P_{jj}}$.
Roughly 99.7\% of errors should lie inside the envelope, and the envelope should not be
grossly loose.

\paragraph{RMSE versus reported $\sigma$} across Monte Carlo runs.

\subsection{IEKF-Specific Diagnostics}
\begin{itemize}
  \item \textbf{Iterations per update.} One to three is typical. Frequently hitting $i_{\max}$
        signals a poor prior, strong nonlinearity, or a bad measurement.
  \item \textbf{Cost monotonicity.} $J(\x\iter{i})$ should decrease at every iteration;
        increases indicate that damping is needed.
  \item \textbf{EKF-versus-IEKF update gap}, $\|\x\iter{N} - \x\iter{1}\|$ in the Mahalanobis
        metric. A large gap flags updates where linearization mattered and where the
        covariance is least trustworthy.
  \item \textbf{Post-fit residuals} $\y_k - \meas(\xhat_k)$ at the converged iterate, compared
        against $\R_k$. A post-fit/pre-fit ratio near unity for every measurement suggests the
        filter is not learning from the data.
\end{itemize}

\subsection{Numerical and Structural Health}
\begin{itemize}
  \item \textbf{Covariance sanity:} symmetry ($\|\Pm - \Pm^\top\|$), positive definiteness
        (Cholesky succeeds; $\lambda_{\min} > 0$), and condition number. Consider a square-root
        or UD formulation if conditioning degrades.
  \item \textbf{Observability:} rank and conditioning of the observability Gramian or
        accumulated information matrix. A collapsing $\sigma$ along an unobservable direction is
        a red flag.
  \item \textbf{Outlier gating:} reject or down-weight measurements with
        $\epsilon_{\nu,k} > \chi^2_m(1 - \alpha)$, and log the rejection rate, which is itself a
        diagnostic.
\end{itemize}

%======================================================================
\section{Use on Real Data}

The IEKF is widely used operationally, including in orbit determination (range, angles,
Doppler), star-tracker/gyro attitude estimation, GNSS/INS integration, radar and
bearings-only tracking, and visual-inertial odometry and SLAM, where many back ends are
effectively iterated EKFs on a manifold. Real data does, however, change how the filter is
validated and tuned.

\begin{itemize}
  \item \textbf{No truth, hence no NEES.} The primary health checks become the
        innovation-based tests: NIS, whiteness, bias, and post-fit residuals. Supplement these
        with an independent reference where available (precise ephemerides, a higher-grade
        sensor, survey points).
  \item \textbf{Validate in simulation first.} Tune the filter and pass NEES/NIS tests in Monte
        Carlo using a realistic truth model, including biases, colored noise, and model
        mismatch. Then replay real data and confirm the NIS statistics still hold.
  \item \textbf{Account for real-world error sources} the textbook model omits:
  \begin{itemize}
    \item sensor biases and scale factors (augment the state);
    \item time-tag misalignment;
    \item colored or heavy-tailed noise (shaping filters, or robust Huber weighting in $J$);
    \item outliers (gating);
    \item unmodeled dynamics (tune $\Q$, or use consider parameters).
  \end{itemize}
  \item \textbf{Tune $\Q$ and $\R$ using NIS.} If the innovation statistics are consistent and
        white, the filter's reported uncertainty is credible---the practical goal on real data.
  \item \textbf{Cross-check against a batch or smoother solution.} A Gauss--Newton batch
        least-squares fit or a smoother over the same data localizes where the filter's
        assumptions break down.
  \item \textbf{Respect the compute budget.} In real-time use, cap $i_{\max}$ and use damping
        so that a pathological measurement degrades gracefully to roughly EKF behavior rather
        than stalling or diverging.
\end{itemize}

%======================================================================
\nocite{sarkka2013}
\printbibliography

\end{document}
